% Expository recovery from the owners of the TPK core.
% Provenance and preservation control are documented in technical/TPK_INTEGRACION.md.
% This file expands the TPK section; it does not replace extension.tex or generacion.tex.

\subsubsection{Dynamical state, observables, and dependence of the operations}
\label{tpk-int:estado}

The topological phase kernel organizes the evolution of the two APP sheets
with the orientation determined by TRIT. Its definition comprises the state
on which it acts, the transition rules, and the relations that make
extensions compatible. Equation \eqref{eq:actualizacion} acquires
operational content by specifying which coordinates each of its factors receives
and modifies. Position belongs to the toroidal support; phase
belongs to the calendar; orientation determines transport; the
quotients record the arithmetic evaluation prior to its reduction.
These coordinates are related through the transition and retain distinct
mathematical functions.

The observable projection of a state is
\begin{equation}
 q=(x_+,d_+,x_\times,d_\times,\theta)
 \in\mathcal Q_{\rm obs}
 :=(X_9\times\mathcal D)^2\times\Theta_{108},
 \qquad \mathcal D=\{N,E,S,O\},\quad
 \Theta_{108}=\ZZ/108\ZZ.
 \label{tpk-int:qobs}
\end{equation}
The subscripts distinguish the additive and multiplicative cursors. In the
fibers of complete states over \(q\), the ordered path, the integer evaluations
of both sheets, their residues and quotients, and the orientation and boundary
records are preserved. A precision extension adds the
prefix and the compatible cylinder. Thus, two states with the same
projection \(q\) may preserve different records.

The dependence of the operations has three levels. At the local level,
the selector activates a sheet, transport modifies its position, and the
update calculates the new arithmetic coordinates. At the path level,
these transitions compose and preserve the records of
preceding steps. At the refinement level, extension relations
determine which continuation preserves the prefix, signature, and
boundary of the antecedent. The finite calendar describes the phase of those
operations; holonomy describes their composition on the history.

\subsubsection{Tritic selection and lifted cardinal transport}
\label{tpk-int:seleccion}

Let \(\bar\theta\in\{0,\ldots,107\}\) be the calendar representative.
The operative phase and the dodecaphase sector are
\begin{equation}
 r(\theta)=1+(\bar\theta\bmod9),\qquad
 \beta(\theta)=\left[\left\lfloor\frac{\bar\theta}{9}\right\rfloor\right]_{12}.
 \label{tpk-int:fase}
\end{equation}
The selector already defined can be written \(\varepsilon_9=M_{\rm ph}\).
Its value vector, in the ordered chart of \(D_9\), is
\[
 m_{\rm ph}=(1,-1,0,1,-1,0,1,-1,0).
\]
The function \(M_{\rm ph}:D_9\to\TRIT\) and this vector are two
presentations of the same selector, with their respective domains.

Define the activation indicators
\[
 a_+(\theta)=\mathbf1_{\{1,4,7\}}(r(\theta)),\qquad
 a_\times(\theta)=\mathbf1_{\{2,5,8\}}(r(\theta)).
\]
The cursor transport is then
\begin{equation}
 x_+'=x_++a_+(\theta)\nu_{d_+},\qquad
 x_\times'=x_\times+a_\times(\theta)\nu_{d_\times}
 \quad\text{in }X_9.
 \label{tpk-int:cursores}
\end{equation}
At phases \(3,6,9\), both positions remain fixed; the neutral event
forms part of the chronology and its record. The cardinal directions
determine the vectors \(\nu_d\), while the phase signature determines
which of those vectors is applied.

The integer lift preserves the boundary crossing. If
\(s:X_9\to D_9^2\) selects the positive representatives and
\(x'=x+a\nu_d\), with \(a\in\{0,1\}\), define
\begin{equation}
 b_d(x,a)=\frac{s(x)+a\nu_d-s(x')}{9}\in\ZZ^2.
 \label{tpk-int:borde}
\end{equation}
Divisibility by \(9\) follows from equality in \(X_9\).
If \(z\in\ZZ^2\) records the preceding crossings, the update is
\(z'=z+b_d(x,a)\). Thus,
\[
 s(x')+9z'=s(x)+9z+a\nu_d.
\]
The integer transport datum remains reconstructible after the
position on the torus is reduced.

\begin{proposition}[Preservation of the lifted displacement]
\label{tpk-int:prop-desplazamiento}
For a cardinal path of \(n\) steps, whose activation indicators
are \(a_t\),
\[
 s(x_n)+9z_n=s(x_0)+9z_0+
 \sum_{t=0}^{n-1}a_t\nu_{d_t}.
\]
The same equality is preserved when two compatible paths are concatenated.
\end{proposition}
\begin{proof}
The one-step identity is \eqref{tpk-int:borde}. Summing it over the
\(n\) steps cancels the intermediate coordinates. Dividing the summation
range into two segments produces exactly the concatenation law. For a
closed path, the difference \(z_n-z_0\) is the previously defined
winding number.
\end{proof}

The same principle governs the APP evaluations. If a sheet takes the
integer value \(N_t\), its record simultaneously preserves
\(N_t=R_t+9q_t\). The difference between two steps satisfies
\[
 N_{t+1}-N_t=(R_{t+1}-R_t)+9(q_{t+1}-q_t).
\]
The spatial-crossing memory \(z_t\) and the arithmetic quotient \(q_t\)
have different definitions; their joint preservation maintains the
geometric and arithmetic provenance of each transition.

\subsubsection{Two-cursor chronology and composite returns}
\label{tpk-int:cronologia}

The cardinal involution \(\jmath\) exchanges \(N\) with \(S\) and
\(E\) with \(O\). After applying \eqref{tpk-int:cursores}, both
directions are reversed when a block of \(27\) steps ends.
Writing
\[
 h(\theta)=\mathbf1_{\{0,27,54,81\}}
             ((\bar\theta+1)\bmod108),
\]
the complete rule on the observable projection is
\begin{equation}
 F_{\rm obs}(q)=
 \bigl(x_+',\jmath^{h(\theta)}d_+,
       x_\times',\jmath^{h(\theta)}d_\times,
       \theta+1\bigr).
 \label{tpk-int:fobs}
\end{equation}
The reading that forms the emission uses the positions prior to
transport, according to the explicit recurrence of
Section~\ref{sec:extension}. The observable update and the emission
therefore compose in a fixed order.

\begin{proposition}[Reversibility and period of the observable calendar]
\label{tpk-int:periodo}
The map \(F_{\rm obs}\) is a permutation of
\(\mathcal Q_{\rm obs}\) and satisfies \(F_{\rm obs}^{108}=I\).
In a block of \(54\) steps, positions and directions are restored;
the \(\Theta_{108}\) coordinate advances by \(54\).
\end{proposition}
\begin{proof}
Given the final state, the preceding phase is first recovered by subtracting one
in \(\Theta_{108}\). That phase determines whether the directions were
reversed and which cursor was displaced. Applying the corresponding involution
and subtracting the cardinal vector uniquely recovers the initial state.

Each interval of \(27\) steps contains nine activations of each
cursor. The following block preserves the activations and reverses the
directions, so the integer displacements cancel over
\(54\) steps. The rule has period \(54\) in its position and
orientation increments; the same cancellation holds for any calendar
origin. After \(108\) steps, \(\theta\) is also restored.
\end{proof}

Set aside
\[
 \mathcal K_{27}=F_{\rm obs}^{27}:
 \mathcal Q_{\rm obs}\longrightarrow\mathcal Q_{\rm obs}
\]
for the iterate of \(27\) transitions. Its fourth iterate is the
identity on the observable state defined in \eqref{tpk-int:qobs}.
Projection of the history to the calendar permits this finite
verification. The path record, its evaluations, and
the new prefixes are preserved in the dynamical state on which
\(\mathcal U_t\) acts; observable equality does not reinitialize them.

\subsubsection{Affine transport of memory and composition of states}
\label{tpk-int:memoria}

The observable projection admits fibers of arithmetic and geometric data.
Over a configuration \(q\), write a state in the form
\[
 x=(q,o,h,c,\sigma,C,b,\lambda).
\]
Here \(o\) preserves the orientation and the sheet; \(h\), the residues
and quotients; \(c\), the additive transport memory; \(\sigma\),
the boundary signature; \(C\), the precision cylinder; \(b\), its
boundary label; and \(\lambda\), the recorded path. The
signature is determined by a map
\[
 \operatorname{Sig}_q:
 \mathsf H_q\times\mathsf C_q\times\mathsf{Hist}_q
 \longrightarrow\mathsf S_q,
 \qquad \sigma=\operatorname{Sig}_q(h,c,\lambda).
\]
The symbols \(\mathsf H_q,\mathsf C_q,\mathsf{Hist}_q\), and
\(\mathsf S_q\) denote, respectively, the sets of residue--quotient data,
the abelian group of transport memory, the
recorded paths, and the boundary signatures.
Realizations of this signature are specified componentwise in
the regional and terminal constructions. This dependence prevents treating
the signature as a free parameter once the data that produce it have been fixed.

The elementary residue-and-quotient chart is explicit. For
\(n\in\ZZ\), let
\[
 r=1+((n-1)\bmod9),\qquad u=\frac{n-r}{9}.
\]
Then \(n=r+9u\), with \(r\in D_9\) and \(u\in\ZZ\).
Uniqueness follows because two representatives of \(D_9\) congruent
modulo \(9\) are equal. This chart provides the local integer
reconstruction. The precision towers also incorporate the truncation maps
and the cylinders developed in
\ref{sec:extension} and \ref{sec:generacion}.

Let \(\gamma:q\to q'\) be an observable path. Transport
of the fibers is expressed by maps
\[
 \mathsf H(\gamma):\mathsf H_q\to\mathsf H_{q'},\qquad
 \mathsf C(\gamma):\mathsf C_q\to\mathsf C_{q'},\qquad
 \mathsf{Hist}(\gamma):\mathsf{Hist}_q\to\mathsf{Hist}_{q'}.
\]
The maps respect identities and composition, and
\(\mathsf C(\gamma)\) is a homomorphism of abelian groups.
The increment \(k_{\rm tr}(\gamma)\in\mathsf C_{q'}\) satisfies
\begin{equation}
 \begin{split}
 k_{\rm tr}(\operatorname{id}_q)&=0,\\
 k_{\rm tr}(\gamma_2\circ\gamma_1)
 &=k_{\rm tr}(\gamma_2)+
   \mathsf C(\gamma_2)k_{\rm tr}(\gamma_1).
 \end{split}
 \label{tpk-int:cociclo}
\end{equation}
Here the symbol \(k_{\rm tr}\) denotes the transport cocycle;
the rational coordinate \(\kappa\) of record \(K\) retains
its notation in the chapter on dodecaphase closure.

The transportable coordinates are updated by
\begin{equation}
 \begin{aligned}
 h'&=\mathsf H(\gamma)h,&
 c'&=\mathsf C(\gamma)c+k_{\rm tr}(\gamma),\\
 \lambda'&=\gamma\circ\lambda,&
 \sigma'&=\operatorname{Sig}_{q'}(h',c',\lambda').
 \end{aligned}
 \label{tpk-int:accion-memoria}
\end{equation}
The second equation is an affine action: the linear term transports
the preceding memory, and the cocycle adds the increment produced by the
path. In the spatial-crossing realization of each cursor separately,
\(\mathsf C(\gamma)=I\), and \(k_{\rm tr}(\gamma)\) is the sum
of the vectors in \eqref{tpk-int:borde}. This concrete instance
links the abstract law with the cardinal operations already defined.

\begin{proposition}[Compatibility of the update with concatenation]
\label{tpk-int:composicion-memoria}
Update \eqref{tpk-int:accion-memoria} by
\(\gamma_1\), followed by update by \(\gamma_2\),
coincides with update by \(\gamma_2\circ\gamma_1\).
\end{proposition}
\begin{proof}
For the affine coordinate, two updates give
\[
 c''=\mathsf C(\gamma_2)\mathsf C(\gamma_1)c
       +\mathsf C(\gamma_2)k_{\rm tr}(\gamma_1)
       +k_{\rm tr}(\gamma_2).
\]
Functoriality of \(\mathsf C\) and
\eqref{tpk-int:cociclo} turn this expression into
\(\mathsf C(\gamma_2\circ\gamma_1)c+
k_{\rm tr}(\gamma_2\circ\gamma_1)\).
The assertion for \(h\) follows from composition of
\(\mathsf H\), and that for \(\lambda\) from associativity of
paths. The final signature coincides because it is evaluated on the same
three reconstructed coordinates.
\end{proof}

The cylinder and its boundary complete this law through the relations
\(\operatorname{Ext}_\gamma\), defined in
\ref{tpk-int:bidireccionalidad}. Their elements determine which of
the preceding updates preserve an admissible extension.
The states, together with the triples \((\gamma,x,x')\) admitted by
those relations, form a category over the observable paths:
\[
 (\gamma_2,x',x'')\circ(\gamma_1,x,x')
     =(\gamma_2\circ\gamma_1,x,x'').
\]
The identity of \(x\) is \((\operatorname{id}_q,x,x)\).
Composition simultaneously preserves arithmetic transport and
boundary extendability. This is the structure in which
returns with memory and successive refinements are effected.

\subsubsection{Dodecaphasic record, phase, and accumulated memory}
\label{tpk-int:dodecafase}

The \(108\) transitions are distributed among the ordered windows
\[
 E_m=\{9m-8,\ldots,9m\},\qquad 1\leq m\leq12.
\]
The record preserves the phase origin, orientation, and data
produced during each window. Its map is
\[
 \mathcal R_{12}(x)=
 \bigl((E_m,\tau_m,\sigma_m,\kappa_m,\Lambda_m)\bigr)_{m=1}^{12},
\]
where \(\tau_m\) is the subroute in the window, \(\sigma_m\) its
orientation, \(\kappa_m\) the integer boundary increment, and
\(\Lambda_m\) the local incidence record. We adopt the chart
to be developed in \ref{subsec:k-registro-integral}; its window indices
correspond to \(m=\beta+1\).
The domain comprises the compatible TPK histories, and the codomain their
oriented temporal records. Composition of segments and their increments
is governed by \eqref{tpk-int:accion-memoria}. The group \(C_{12}\), by
contrast, describes translation of the window index; \(\Theta_{108}\) preserves
the step calendar. This distinction specifies what is transformed
and what is observed upon completion of a cycle.

In coordinates \(\theta=9b+r\), with \(0\leq r<9\), advancing
\(9\) steps preserves \(r\) and increases \(b\) by one.
Addition of phases introduces the quotient
\[
 c_0(r,r')=\left\lfloor\frac{r+r'}9\right\rfloor,
\]
whose cocycle identity ensures associativity of composition.
The finite calendar preserves \([b]_{12}\); a record of
circuits preserves \(b\in\ZZ\), or \(b\in\mathbb N_0\) for
a history with fixed origin. Section~\ref{sec:extension} proves
the nonsplit exact sequence \eqref{eq:extension108}. Its function is
to describe the relation between these coordinates before using them in the
construction of the signed state.

The signed record receives the phase and memory data through the
composition developed in Section~\ref{sec:k-registro}; on the
compatible sublattice, the notations \(\mathcal S_{12}\) and
\(R_{\mathrm{sgn}}\) denote the same reader:
\begin{equation}
 \mathcal S_{12}\equiv R_{\mathrm{sgn}}
 =\Pi_H\circ\mathcal E_{108}^{90,120}
                         \circ\mathcal R_{12}.
 \label{tpk-int:registro-causal}
\end{equation}
Its domain is the terminal state with memory; its output belongs to
\(\ZZ^{12}\). The reversible transformation between that signed state
and \(K\), as well as the rational reading of \(K\), are proved in
the corresponding chapter. The recording operations precede
those transformations. Reconstruction of a record from its cyclic
differences verifies its recoverability; the dynamical provenance
must be preserved through the events that produced it.
Equation~\eqref{tpk-int:registro-causal} fixes the type and effective
composition. The terminal state is produced by the one hundred eight updates
\(\mathcal U_t=\mathsf{Upd}_t\circ\mathsf{Tra}_t\circ\mathsf{Sel}_t\);
\(\mathcal R_{12}\) preserves the twelve windows with orientation, carry,
boundary, and ledger, and \(\mathcal E_{108}^{90,120}\) expresses the channels
\((b^{90},b^{120},Q_{\rm TPK})\). Section~\ref{sec:k-registro}
incorporates the exact evaluation and proves that
\(\Pi_H\) produces \(U_{\rm term}\) before the Hadamard inversion that
produces \(K\). The focal verifier deliberately begins at the channel
coordinate to check the subsequent chart and its return focally and exactly;
the scope of that auxiliary program is not confused with the scope of
the mathematical construction.

For the coded trace \((d_t)\) preserved by the record, the
event reader counts codes \(\{2,4,6,8\}\) in each
window with orientation
\((+,+,+,-,-,-,+,+,+,-,-,-)\). These codes and the cardinal
symbols in \(\mathcal D\) are distinct coordinates of the
record. Chapter~\ref{sec:k-registro} develops this reader,
the integral decomposition \(1+5+6\), the congruences that
characterize its image, and its exact inverse. The inversion proof
establishes which data that chart preserves; the definition of
\(\mathcal R_{12}\) specifies the oriented history on which it
acts.

\subsubsection{Phase incidence and internal construction of the channels}
\label{tpk-int:incidencia}

The selector determines three subsets of \(D_9\):
\[
 H_+=\{1,4,7\},\qquad H_-=\{2,5,8\},\qquad H_0=\{3,6,9\}.
\]
Let \(H_{\rm act}=H_+\sqcup H_-\) be the set of active phases and
let \(O_{\rm ph}=D_9\setminus\{9\}\) be the chart preceding the marked
return. Denote by \(P_2(H_{\rm act})\) the set of unordered pairs
of active phases. The internal incidence is given by
\begin{equation}
 \mathcal F^{\rm ph}_6=H_{\rm act}\times P_2(H_{\rm act}),\qquad
 \mathcal F^{\rm ph}_8=O_{\rm ph}\times P_2(H_{\rm act}),\qquad
 \Theta_{\rm ph}=D_9\times P_2(H_{\rm act}).
 \label{tpk-int:fases90120}
\end{equation}
The indices describe the number of positions in the first component.
The construction uses the phases and calendar origin already fixed;
the subsequent exceptional realizations will transport these
incidences to their own domains.

\begin{proposition}[Cardinalities and oriented phase incidence]
\label{tpk-int:cardinales}
The sets in \eqref{tpk-int:fases90120} have cardinalities
\(90,120,135\), respectively. If
\[
 \partial_{\rm ph}=\Theta_{\rm ph}\times\{-1,+1\},\qquad
 E_{\rm ph}=\partial_{\rm ph}\sqcup\{\infty\},\qquad
 L_{\rm ph}=H_0\times P_2(H_{\rm act}),
\]
then
\[
 |\partial_{\rm ph}|=270,\quad |E_{\rm ph}|=271,\quad
 |L_{\rm ph}|=45,\quad
 |\partial_{\rm ph}\sqcup L_{\rm ph}|=315.
\]
Simultaneously translating the phase origin and the marked return preserves
all these cardinalities and their inclusions.
\end{proposition}
\begin{proof}
There are six active phases, so
\(|P_2(H_{\rm act})|=\binom62=15\). Multiplication by
\(6,8,9\) gives \(90,120,135\). Orientation doubles \(135\);
the return mark adds one element; the three neutral phases give
\(3\cdot15=45\). A calendar translation transports each
factor by a bijection and takes the marked point to its image. Products
and disjoint unions thus preserve incidence.
\end{proof}

Here the cardinality \(271\) has a realization by phase incidence.
The cubic completion of Section~\ref{sec:extension} has
its own construction of the shell \(1000-729\). The correspondence
between the two realizations must preserve the maps and marks that
identify them, in addition to the cardinality. Likewise, the selector
\(\varepsilon_9\), the incidence \((90,120)\), and the extractor
\(\mathcal E_{108}^{90,120}\) preserve the domains and operations
just distinguished.

\subsubsection{Path reversal and conjugate extensions}
\label{tpk-int:bidireccionalidad}

Bidirectionality begins with an operation on paths, before
scalar values are assigned to them. For
\(\gamma=(x_0;d_1,\ldots,d_n)\), with endpoint \(x_n\), define
\[
 \operatorname{rev}\gamma
   =(x_n;\jmath d_n,\ldots,\jmath d_1).
\]
The origin of the reversed path is the endpoint of the original.
Composition satisfies
\[
 \operatorname{rev}^2\gamma=\gamma,\qquad
 \operatorname{rev}(\gamma_2\circ\gamma_1)
   =\operatorname{rev}\gamma_1\circ\operatorname{rev}\gamma_2,
\]
and the integer displacement changes sign. These identities are obtained
by reversing the order of the steps and replacing each cardinal vector by
its opposite. For a loop it follows that
\(\operatorname{wind}(\operatorname{rev}\gamma)
=-\operatorname{wind}(\gamma)\).

On the records, reversal also requires transporting the order
of their components and the orientation of their events. Route reversal
and signature conjugation are distinguishable operations: their
composition is specified when applied to a particular record. In
particular, the regional involution that exchanges the closure and
propagation components preserves the self-scaling axis; its formula is
presented in Section~\ref{mono-ampl:regiones}. This regional operation
acts on a domain different from reversal of a cardinal path.

Extensions along a route \(r:q\to q'\) are expressed by
relations between fibers of compatible states. If \(\mathcal F_q\)
denotes the fiber of complete states over \(q\), require
\[
 \operatorname{Ext}_r\subseteq\mathcal F_q\times\mathcal F_{q'},
 \qquad \Delta_{\mathcal F_q}\subseteq
        \operatorname{Ext}_{\operatorname{id}_q},
\]
\[
 \operatorname{Ext}_{r_2}\circ\operatorname{Ext}_{r_1}
       \subseteq\operatorname{Ext}_{r_2\circ r_1}.
\]
Their composition preserves the evaluations of both sheets, the orientation,
the quotients, the prefix, and the boundary. Reversal of the route does not
by itself imply that every extension relation is a bijection.
The conjugate extension is defined on the domain where transport of
those data satisfies the corresponding compatibilities.

\subsubsection{Nonadic holonomy and continuity of refinement}
\label{tpk-int:holonomia}

One phase circuit composes the transitions on the state with memory:
\begin{equation}
 \Gamma_{9,k}=\mathcal U_{k+8}\circ\cdots\circ\mathcal U_k.
 \label{tpk-int:gamma}
\end{equation}
On an individualized history, this functional composition realizes
the holonomy relation; on the complete domain,
\(\Gamma_{9,k}\subseteq\widehat X_k\times\widehat X_{k+9}\) is preserved,
with admissible extensions developed below.
The cyclic base has nine phases, while
the fiber preserves the path, quotients, and boundary.
Monodromy is the action of one circuit of that base on the fibers; the
holonomy \(\Gamma_{9,k}\) expresses its realization by the effective
TPK transitions.

If \(g\) observes the phase and \(M\) the number of preserved circuits,
\[
 g(\Gamma_{9,k}x)=g(x),\qquad M(\Gamma_{9,k}x)=M(x)+1.
\]
The first identity expresses phase closure, and the second identifies
the new temporal depth. Coinductive extension also preserves
its cylinders and prefixes. The restrictions between depths
compose an inverse system; a compatible section gathers its
antecedents without replacing them by the latest observation. The development
of this arithmetic of histories and its identity occupies
\ref{hist:seccion}--\ref{hist:doble-varianza}; construction of the
regions and their readers is carried out in Section~\ref{sec:generacion}.

Phase transfer on already constructed emissions has a different
domain. Let \(\mathcal A_+\) and \(\mathcal A_\times\) be the
alphabets of additive and multiplicative signatures, of cardinalities \(18\)
and \(26\). Their product identifies the catalog \(\mathcal U_6\).
For \(f:\mathcal A_+\times\mathcal A_\times\to\RR\),
the sheet averages are
\[
 (E_+f)(s,e)=\frac1{18}\sum_{s'\in\mathcal A_+}f(s',e),\qquad
 (E_\times f)(s,e)=\frac1{26}\sum_{e'\in\mathcal A_\times}f(s,e').
\]
These emissions and their cardinalities are constructed by means of
the recurrence in Section~\ref{sec:extension}. Once obtained,
the transfer operator is defined on
\(\mathcal U_6\times C_9\) by
\[
 P_+=E_+,\quad P_-=E_\times,\quad P_0=I,\qquad
 \mathcal K_{\rm ph}((u,r),(v,r+1))=P_{M_{\rm ph}(r)}(u,v),
\]
with zero weight for the remaining phase transitions. The index
\(r\) uses the positive representative of the cyclic class.
The phase order \((+1,-1,0)^3\) then produces the renewal
\[
 \mathcal K_{\rm ph}^{9}\big|_r=E_+E_\times.
\]
Indeed, both averages are idempotent, commute, and average different
factors; their composition assigns weight \(1/(18\cdot26)=1/468\)
to each emission. This equality characterizes a subsequent transfer
observable. Transport of the histories remains given by
\eqref{tpk-int:gamma}, with its memory and boundary.

The resulting hierarchy determines the causal order of subsequent
constructions.
The finite emission constructs the domain and regions; extension relations
preserve their prefixes; the nonadic connection extends them; the
dodecaphase register determines the data entering \(K\) and the
closure of \(\alpha\). The numerical and incidence realizations
receive that preceding structure. This order makes it possible to present
each consequence in its chapter while preserving, from the formal core onward,
the objects that make it possible.
